%% ma: L12-B08-0090
%% lop: 12
%% chuong: 2
%% bai: 8
%% dang: 12.8.12
%% loai: TN4
%% muc: TH
%% dap_an: "B"
%% nguon: "CĐ2.2.pdf giáo viên gửi 27/10/2026 (M07, Lê Minh Khôi – Nguyễn Anh Khôi, Chương II Chủ đề 2 Tọa độ của vectơ trong không gian), Câu 24"
%% kien_thuc: "Bài 8 (tọa độ vectơ, hai vectơ bằng nhau), Bài 6 (vectơ bằng nhau trong hình hộp)"
%% trang_thai: chinh-thuc
\begin{ex}
Trong không gian $Oxyz$, cho hình hộp chữ nhật $ABCD.A'B'C'D'$ có tọa độ các đỉnh $B'(2;-4;0)$, $C'(-1;1;-1)$, $A(3;5;1)$ và $C(-2;0;4)$. Tọa độ của điểm $D'$ là
\choice{$D'(11;1;-5)$}{\True $D'(1;11;-5)$}{$D'(-5;1;11)$}{$D'(1;-5;11)$}
\loigiai{$\overrightarrow{AA'}=\overrightarrow{CC'}=(1;1;-5)\Rightarrow A'(4;6;-4)$.

$\overrightarrow{A'D'}=\overrightarrow{B'C'}=(-3;5;-1)\Rightarrow D'(1;11;-5)$.}
\end{ex}
